% ----------------------------------------------------------
% Desenvolvimento
% ----------------------------------------------------------
\chapter{Desenvolvimento da solução}
\label{cap:desenvolvimento}
% ----------------------------------------------------------

Este capítulo descreve a solução implementada no repositório. A apresentação procura refletir o estado atual do sistema, sem tratar como concluídas etapas que ainda estão em validação. A aplicação foi construída como um sistema web de consulta a manuais automotivos, com processamento local de documentos, recuperação textual e visual, memória conversacional e geração de respostas com modelo local.

\section{Arquitetura geral}

A arquitetura é composta por três partes principais. A primeira é o frontend, desenvolvido em Next.js, responsável pela interface de chat, exibição de fontes, confiança e imagens. A segunda é o backend, desenvolvido em NestJS, responsável por autenticação, chats, estado conversacional, orquestração do RAG, geração e pós-processamento. A terceira é o conjunto de scripts Python, responsável por consultar o índice FAISS, aplicar sinais lexicais e analisar imagens com CLIP.

\begin{figure}[H]
\centering
\begin{minipage}{0.88\textwidth}
\begin{verbatim}
Usuário
  -> Frontend Next.js
  -> Backend NestJS
  -> Conversation State e Query Rewriter
  -> Recuperação Python
  -> FAISS + BM25 + regras de domínio
  -> Ollama / DeepSeek R1 8B
  -> Pós-processamento, confiança, fontes e imagens aprovadas
\end{verbatim}
\end{minipage}
\caption{Fluxo lógico da solução implementada.}
\fonte{Elaborado pelo autor.}
\label{fig:arquitetura}
\end{figure}

O sistema mantém os documentos processados em uma estrutura local chamada \texttt{rag\_chunks\_output}, contendo o índice FAISS, os metadados dos chunks e as imagens extraídas. O banco relacional, acessado via Prisma, é utilizado para usuários, chats, mensagens e metadados conversacionais.

\section{Frontend}

O frontend permite criar conversas, enviar perguntas, receber respostas por streaming e visualizar fontes utilizadas. As mensagens do assistente podem trazer documento, páginas, seção, score, confiança, consulta resolvida, status e imagens. Uma alteração importante realizada durante o desenvolvimento foi impedir que a interface exibisse imagens candidatas não aprovadas. Antes, a galeria podia utilizar caminhos de imagens associados aos chunks, mesmo quando o backend não havia confirmado a relevância visual. A versão atual passa a renderizar apenas imagens aprovadas nos metadados da resposta ou em \texttt{clip\_results}.

Essa decisão é relevante para a qualidade do sistema. Em perguntas visuais, uma imagem incorreta pode prejudicar mais do que a ausência de imagem. Por isso, o frontend passou a respeitar o contrato do backend: se nenhuma imagem for validada, a resposta permanece textual.

\section{Backend}

O backend concentra a lógica conversacional e a orquestração do RAG. A camada de chats recupera o histórico recente e o último estado conversacional salvo. Em seguida, o serviço RAG executa sanitização, reescrita da consulta, detecção de intenção, recuperação, geração, validação e persistência da resposta.

As principais etapas do backend são:

\begin{enumerate}
  \item sanitizar a mensagem e o histórico;
  \item carregar o estado conversacional anterior;
  \item identificar se a mensagem é pergunta nova, complemento ou continuação;
  \item montar uma consulta resolvida;
  \item recuperar chunks relevantes;
  \item gerar resposta com o modelo local;
  \item validar resposta, fontes, confiança e imagens;
  \item salvar metadados da resposta para uso nas próximas mensagens.
\end{enumerate}

Esse fluxo foi necessário porque mensagens curtas, como ``2008'', ``2010'' ou ``2.0'', não possuem sentido suficiente quando analisadas isoladamente. No sistema atual, essas mensagens podem completar uma pergunta pendente antes da busca documental.

\section{Processamento documental}

O processamento dos PDFs é realizado em um notebook Python. O pipeline extrai texto, imagens, páginas, seções e metadados. Em seguida, cria chunks semânticos com base em títulos, limites de página e tamanho do texto. Cada chunk preserva informações como documento, páginas, seção, texto, imagens associadas e metadados do veículo quando inferidos.

Os embeddings são gerados com Sentence Transformers e armazenados em um índice FAISS. Como os vetores são normalizados, a busca por produto interno se aproxima da similaridade de cosseno. O arquivo \texttt{chunks\_metadata.json} mantém a ligação entre cada vetor e seu conteúdo textual, documento, páginas e imagens.

\section{Recuperação híbrida}

O script \texttt{retrieve\_chunks.py} implementa a recuperação documental. A busca não depende apenas da similaridade vetorial. Ela combina:

\begin{itemize}
  \item busca vetorial com FAISS;
  \item BM25 local;
  \item pontuação de domínio para óleo, pneus, rodas e localização visual;
  \item compatibilidade com ano, marca e modelo;
  \item reranker textual opcional por \textit{cross-encoder}.
\end{itemize}

Essa combinação foi incorporada porque perguntas automotivas exigem precisão documental. Em uma pergunta sobre a tampa de óleo, por exemplo, não basta recuperar qualquer trecho sobre óleo. O sistema deve priorizar trechos com ``engine oil filler cap'', ``oil filler cap'', ``engine compartment'' ou figuras de localização. Também foi reduzido o bônus dado a qualquer chunk com imagem, pois imagens próximas no PDF nem sempre representam a peça perguntada.

\section{Memória conversacional e reescrita}

O estado conversacional é salvo nos metadados das mensagens do assistente. Ele inclui a última pergunta do usuário, pergunta pendente, consulta resolvida, intenção, veículo, slots ausentes e indicador de esclarecimento pendente. O módulo \texttt{QueryRewriterService} usa esse estado para transformar a mensagem atual em uma consulta completa.

Um exemplo prático é a sequência:

\begin{verbatim}
Usuário: Onde fica a tampa do óleo do Audi A4?
Sistema: Informe o ano do veículo.
Usuário: 2008
\end{verbatim}

Nesse caso, a busca não deve receber apenas ``2008''. O backend monta uma consulta resolvida como ``Onde fica a tampa do óleo do Audi A4 2008?''. Essa etapa melhora a continuidade conversacional e evita respostas deslocadas da intenção original.

\section{Geração local de respostas}

A geração textual é realizada com Ollama e DeepSeek R1 8B. O prompt instrui o modelo a responder em português, usar apenas o contexto recuperado, evitar especulações e citar fontes. Apesar disso, durante os testes observou-se que o modelo ainda podia gerar respostas pouco úteis ou confundir termos, como tratar a tampa de óleo como um reservatório. Por esse motivo, foi adicionada uma camada de pós-processamento no backend.

Essa camada verifica se a resposta apenas manda consultar o manual, se usa termos proibidos para o contexto, se está genérica ou se contradiz a evidência recuperada. Para perguntas de localização de componentes com evidência clara, como ``engine oil filler cap'', o backend pode montar uma resposta determinística a partir do manual. Assim, a geração do modelo deixa de ser a única responsável pela resposta final.

\section{Pipeline de imagens}

O pipeline visual utiliza imagens extraídas dos PDFs durante a ingestão. O CLIP é acionado somente quando a consulta resolvida possui intenção visual e existem imagens candidatas. O script \texttt{analyze\_images\_with\_clip.py} calcula similaridade entre imagens e prompts textuais. A versão atual utiliza prompts positivos e negativos. Para uma pergunta sobre tampa de óleo, por exemplo, prompts positivos incluem diagramas do compartimento do motor com reservatórios e tampa de óleo; prompts negativos incluem alavanca de abertura do capô, trava do capô e imagens externas do veículo.

Depois do CLIP, o backend executa um reranqueamento adicional em \texttt{VisualImageService}. Essa etapa considera:

\begin{itemize}
  \item score positivo do CLIP;
  \item score negativo do CLIP;
  \item legenda extraída do texto do chunk;
  \item OCR ou texto ao redor da figura;
  \item presença de termos da peça perguntada;
  \item rejeição de capa, imagem vazia, imagem genérica e figura sem relação com a pergunta.
\end{itemize}

Esse mecanismo foi motivado por falhas reais. Em um caso, a pergunta era sobre a tampa de óleo, mas a imagem exibida mostrava a abertura do capô. A correção implementada passou a rejeitar imagens que o CLIP identifica como alavanca ou trava do capô quando a pergunta é sobre tampa de óleo. Também foi validado que, para o Audi A4 2008, a recuperação textual prioriza o trecho da Figura 252, que apresenta o layout do compartimento do motor e a indicação da tampa de abastecimento do óleo.

\section{Confiança, fontes e logs}

A resposta final inclui score de confiança, fontes, páginas, citações, imagens aprovadas e avisos. O cálculo de confiança considera scores de recuperação, metadados, consenso de documentos, qualidade das fontes e resultado da validação visual. Penalizações são aplicadas quando a resposta é genérica, quando apenas orienta consultar o manual, quando a fonte não sustenta a resposta ou quando imagens são rejeitadas.

O backend também registra logs estruturados com mensagem crua, pergunta pendente, consulta resolvida, páginas recuperadas, imagens selecionadas, imagens rejeitadas, razões de rejeição e confiança antes e depois das penalizações. Esses logs são importantes para depurar conversas reais e compreender por que uma resposta foi montada de determinada forma.

\section{Estado atual da implementação}

Até o momento, estão implementados: aplicação web, autenticação, chats, persistência de mensagens, pipeline de ingestão, FAISS, recuperação híbrida, geração local, memória conversacional, reescrita de consultas, CLIP local, reranking de imagens, bloqueio de imagens não aprovadas no frontend, cálculo de confiança e testes automatizados para casos críticos.

Ainda estão em andamento: calibração dos pesos de confiança, ampliação do conjunto de testes com manuais reais, melhoria da associação entre figura e legenda, avaliação quantitativa com modelos externos e construção de um benchmark automotivo. Essas etapas fazem parte da continuidade natural do trabalho e serão necessárias para consolidar a análise comparativa proposta no título.
